\documentclass[12pt,a4paper]{article}

\usepackage[margin=1in]{geometry}
\usepackage[hidelinks]{hyperref}
\usepackage{setspace}
\usepackage[ngerman]{babel}
\usepackage{changepage}
\usepackage{newtxtext}

\pagestyle{empty}

\begin{document}

\begin{center}
{\Huge Die Geburt der Tragödie}\\
{\Large Quin Barnard}
\end{center}

\hfill den 27. der 8. 2026

\textbf{\large Friedrich Nietzsche}
\begin{itemize}
    \item 15. Oktober 1844 bis 25 August 1900
    \item Fängt im 1864 an Theologie zu studieren; Bonn
    \item Wechselt im 1865 um: Klassische Philologie, Leipzig
    \item Veröffentlicht 1872 \textit{„Die Geburt der Tragödie“}
\end{itemize}

\noindent
\begin{minipage}[t]{0.48\textwidth}

    \textbf{\Large{Apollon}}
        \begin{itemize}
            \item Dichtkunst
            \item Vernunft
            \item Sichtbar
            \item Form
            \item Grenzen
            \item \textbf{Principium Individuationis}
        \end{itemize}

\end{minipage}
\hfill
\begin{minipage}[t]{0.48\textwidth}

    \textbf{\Large Dionysos}
        \begin{itemize}
            \item Theater
            \item Rausch
            \item Unsichtbar
            \item Formlosigkeit
            \item Grenzlos
            \item \textbf{Zerstört Principium Individuationis}
        \end{itemize}

\end{minipage}

\medskip\noindent\textbf{\Large Die Geburt der Tragödie}
\begin{itemize}    
    \item Andere Theorien des Apollinischen; Johann Joachim Winckelmann
    \item Nietzsche: Der Mensch beinhaltet sowohl das Apollinische als auch das Dionysische
    \item Streit zwischen den Beiden=die Altgriechische Tragödie
    \item Muss Beide haben: \begin{itemize}
                                \item ohne Apollon—Vergessene selbst; ohne Dionysos—Kein Selbst zu definieren übrig
                            \end{itemize}
\end{itemize}

\medskip\noindent\textbf{\Large Quellen}\\
{\small Groddeck, Wolfram: “Die Geburt der Tragödie” in “Ecce homo”: Hinweise zu einer strukturalen Lektüre von Nietzsches Ecce homo. \textit{Nietzsche–Studien. Internationales Jahrbuch für die Nietsche–Forschung} 1984(13):325–331\\Nietzsche, Friedrich: Die Geburt der Tragödie/Aus Dem Gedankenkreise der Geburt der Tragödie. In: Friedrich Niezsche Gesammelte Werke. Dritter Band, Musarion Verlag, München 1920.\\Schlögl, Albert: Friedrich Nietzsches Theorie zur Entstehung der griechischen Tragödie. In: Bohnstingl, Ren\'{e}; Kranebitter, Andreas; Obermayr, Linda Lilith und Reitter, Karl (Hg.): Jahrbuch für Marxistische Gesellschaftstheorie. Berlin 2022, S.142-155}

\end{document}
